% ==== VIII ====
\yoschapter{VIII}{Resolvent and Spectrum}
\label{chap:VIII}

Let $T$ be a linear operator whose domain $D(T)$ and range $R(T)$ both lie in the same complex linear topological space $X$. We consider the linear operator 
\[
T _{\lambda} = \lambda I - T,
\]
where $\lambda$ is a complex number and $I$ the identity operator. The distribution of the values of $\lambda$ for which $T_{\lambda}$ has an inverse and the properties of the inverse when it exists, are called the spectral theory for the operator $T$. We shall thus discuss the general theory of the inverse of $T_{\lambda}$. 


\yossection{1}{The Resolvent and Spectrum}
\label{sec:VIII.1}

\begin{definitionu} If $\lambda_0$ is such that the range $R(T_{\lambda_0})$ is dense in $X$ and $T_{\lambda_0}$ has a continuous inverse $(\lambda_0I - T)^{-1}$, we say that $\lambda_0$ is in the resolvent set $\varrho(T)$ of $T$, and we denote this inverse $(\lambda_0I - T)^{-1}$ by $R(\lambda_0; T)$ and call it the resolvent (at $\lambda_0$) of $T$. All complex numbers $\lambda$ not in $\varrho(T)$ form a set $\sigma(T)$ called the spectrum of $T$. The spectrum $\sigma(T)$ is decomposed into disjoint sets $P_{\sigma}(T)$, $C_{\sigma}(T)$ and $R_{\sigma}(T)$ with the following properties: $P_{\sigma}(T)$ is the totality of complex numbers $\lambda$ for which $T_{\lambda}$ does not have an inverse; $P_{\sigma}(T)$ is called the point spectrum of $T$. 

$C_{\sigma}(T)$ is the totality of complex numbers $\lambda$ for which $T_{\lambda}$ has a discontinuous inverse with domain dense in $X$; $C_{\sigma}(T)$ is called the continuous spectrum of $T$. 

$R_{\sigma}(T)$ is the totality of complex numbers $\lambda$ for which $T_{\lambda}$ has an inverse whose domain is not dense in $X$; $R_{\sigma}(T)$ is called the residual spectrum of $T$. 
\end{definitionu}

From these definitions and the linearity of $T$ we have the 

\begin{propositionu} $A$ necessary and sufficient condition for $\lambda_{0}\in P_{\sigma}(T)$ is that the equation $Tx=\lambda_{0}x$ has a solution $x\neq0$. In this case $\lambda_{0}$ is called an eigenvalue of $T$, and $x$ the corresponding eigenvector. The null space $N(\lambda_{0}I-T)$ of $T_{\lambda_{0}}$ is called the eigenspace of $T$ corresponding to the eigenvalue $\lambda_{0}$ of $T$. It consists of the vector $0$ and the totality of eigenvectors corresponding to $\lambda_{0}$. The dimension of the eigenspace corresponding to $\lambda_{0}$ is called the multiplicity of the eigenvalue $\lambda_{0}$. 
\end{propositionu}

\begin{theoremu} Let $X$ be a complex $B$-space, and $T$ a closed linear operator with its domain $D(T)$ and range $R(T)$ both in $X$. Then, for any $\lambda_{0} \in \varrho(T)$, the resolvent $(\lambda_{0}I - T)^{-1}$ is an everywhere defined continuous linear operator. 
\end{theoremu}

\begin{proof}
 Since $\lambda_0$ is in the resolvent set $\varrho(T), R(\lambda_0 I - T) = D((\lambda_0 I - T)^{-1})$ is dense in $X$ in such a way that there exists a positive constant $c$ for which 
\[
\left\| \left(\lambda_{0} I - T\right) x \right\| \geqq c \| x \| \text{ whenever } x \in D (T).
\]

We have to show that $R(\lambda_0 I - T) = X$. But, if $s\text{- }\lim\limits_{n \to \infty} (\lambda_0 I - T)x_n = y$ exists, then, by the above inequality, $s\text{- }\lim\limits_{n \to \infty} x_n = x$ exists, and so, by the closure property of $T$, we must have $(\lambda_0 I - T)x = y$. Hence, by the assumption that $R(\lambda_0 I - T)^a = X$, we must have $R(\lambda_0 I - T) = X$. 
\end{proof}

\begin{example}\label{ex:VIII.1.Example.1} If the space $X$ is of finite dimension, then any bounded linear operator $T$ is represented by a matrix $(t_{ij})$. It is known that the eigenvalues of $T$ are obtained as the roots of the algebraic equation, the so-called secular or characteristic equation of the matrix $(t_{ij})$: 
\begin{equation}\label{eq:VIII.1.1}
\det \left(\lambda \delta_{ij} - t _{ij}\right) = 0,
\end{equation}
where $\det(A)$ denotes the determinant of the matrix $A$. 
\end{example}

\begin{example}\label{ex:VIII.1.Example.2} Let $X = L_{2}(-\infty, \infty)$ and let $T$ be defined by 
\[
T \cdot x (t) = t x (t),
\]
that is, $D(T) = \{x(t); x(t) \text{ and } tx(t) \in L^2(-\infty, \infty)\}$ and $Tx(t) = tx(t)$ for $x(t) \in D(T)$. Then every real number $\lambda_0$ is in $C_\sigma(T)$. 
\end{example}

\begin{proof}
 The condition $(\lambda_{0}I-T)x=0$ implies $(\lambda_{0}-t)x(t)=0$ a.e., and so $x(t)=0$ a.e. Thus $(\lambda_{0}I-T)^{-1}$ exists. The domain $D((\lambda_{0}I-T)^{-1})$ comprises those $y(t)\in L^{2}(-\infty,\infty)$ which vanish identically in the neighbourhood of $t=\lambda_{0}$; the neighbourhood may vary with $y(t)$. Hence $D((\lambda_{0}I-T)^{-1})$ is dense in $L^{2}(-\infty,\infty)$. It is easy to see that the operator $(\lambda_{0}I-T)^{-1}$ is not bounded on the totality of such $y(t)$'s. 
\end{proof}

\begin{example}\label{ex:VIII.1.Example.3} Let $X$ be the Hilbert space $(l^2)$, and let $T_0$ be defined by 
\[
T _{0} (\xi_{1}, \xi_{2}, \dots) = (0, \xi_{1}, \xi_{2}, \dots).
\]
\end{example}

Then $0$ is in the residual spectrum of $T$, since $R(T_0)$ is not dense in $X$. 

\begin{example}\label{ex:VIII.1.Example.4} Let $H$ be a self-adjoint operator in a Hilbert space $X$. Then the resolvent set $\varrho(H)$ of $H$ comprises all the complex numbers $\lambda$ with $\operatorname{Im}(\lambda) \neq 0$, and the resolvent $R(\lambda; H)$ is a bounded linear operator with the estimate 
\begin{equation}\label{eq:VIII.1.2}
\left\| R (\lambda ; H) \right\| \leqq 1 / \left| \operatorname{Im} (\lambda) \right|.
\end{equation}
\end{example}

Moreover, 
\begin{equation}\label{eq:VIII.1.3}
\operatorname{Im} \left((\lambda I - H) x, x\right) = \operatorname{Im} (\lambda) \| x \| ^{2}, x \in D (H).
\end{equation}

\begin{proof}
 If $x \in D(H)$, then $(Hx, x)$ is real since $(Hx, x) = (x, Hx) = (\overline{Hx}, x)$. Therefore we have \eqref{eq:VIII.1.3}, and so, by Schwarz' inequality, 
\begin{equation}\label{eq:VIII.1.4}
\left\| (\lambda I - H) x \right\| \cdot \left\| x \right\| \geqq \left| ((\lambda I - H) x, x) \right| \geqq \left| \operatorname{Im} (\lambda) \right| \cdot \left\| x \right\| ^{2}
\end{equation}
which implies that 
\begin{equation}\label{eq:VIII.1.5}
\left\| (\lambda I - H) x \right\| \geqq \left| \operatorname{Im} (\lambda) \right| \cdot \left\| x \right\|, x \in D (H). \end{equation}

Hence the inverse $(\lambda I-H)^{-1}$ exists if $\operatorname{Im}(\lambda) \neq 0$. Moreover, the range $R(\lambda I-H)$ is dense in $X$ if $\operatorname{Im}(\lambda) \neq 0$. If otherwise, there would exist a $y \neq 0$ orthogonal to $R(\lambda I-H)$, i.e., $((\lambda I-H)x,y)=0$ for all $x \in D(H)$ and so $(x,(\overline{\lambda}I-H)y)=0$ for all $x \in D(H)$. Since the domain $D(H)$ of a self-adjoint operator $H$ is dense in $X$, we must have $(\bar{\lambda} I - H)y = 0$, that is, $Hy = \bar{\lambda} y$, contrary to the reality of the value $(Hy, y)$. 

Therefore, by the above Theorem, we see that, for any complex number $\lambda$ with $\operatorname{Im}(\lambda) \neq 0$, the resolvent $R(\lambda; H)$ is a bounded linear operator with the estimate \eqref{eq:VIII.1.2}. 
\end{proof}


\yossection{2}{The Resolvent Equation and Spectral Radius}
\label{sec:VIII.2}

\begin{theorem}\label{thm:VIII.2.Theorem.1} Let $T$ be a closed linear operator with domain and range both in a complex $B$-space $X$. Then the resolvent set $\varrho(T)$ is an open set of the complex plane. In each component (the maximal connected sets) of $\varrho(T)$, $R(\lambda; T)$ is a holomorphic function of $\lambda$. 
\end{theorem}

\begin{proof}
 By the Theorem of the preceding section, $R(\lambda; T)$ for $\lambda \in \varrho(T)$ is an everywhere defined continuous operator. Let $\lambda_0 \in \varrho(T)$ and consider 
\begin{equation}\label{eq:VIII.2.1}
S (\lambda) = R \left(\lambda_{0}; T\right) \left\{I + \sum\limits_{n=1} ^{\infty} \left(\lambda_{0} - \lambda\right) ^{n} R \left(\lambda_{0}; T\right) ^{n} \right\}. \end{equation}

The series is convergent in the operator norm whenever $\left|\lambda_{0}-\lambda\right|$. $\left\|R\left(\lambda_{0};T\right)\right\|<1$, and within this circle of the complex plane, the series defines a holomorphic function of $\lambda$. Multiplication by $(\lambda I-T)=(\lambda-\lambda_{0})I+(\lambda_{0}I-T)$ on the left or right gives $I$ so that the series $S(\lambda)$ actually represents the resolvent $R(\lambda;T)$. Thus we have proved that a circular neighbourhood of $\lambda_{0}$ belongs to $\varrho(T)$ and $R(\lambda;T)$ is holomorphic in this neighbourhood. 
\end{proof}

\begin{theorem}\label{thm:VIII.2.Theorem.2} If $\lambda$ and $\mu$ both belong to $\varrho(T)$, and if $R(\lambda; T)$ and $R(\mu; T)$ are everywhere defined continuous operators, then the resolvent equation holds: 
\begin{equation}\label{eq:VIII.2.2}
R (\lambda ; T) - R (\mu ; T) = (\mu - \lambda) R (\lambda ; T) R (\mu ; T).
\end{equation}
\end{theorem}

\begin{proof}
 We have 
\[
\begin{aligned}
&R (\lambda ; T)  = R (\lambda ; T) (\mu I - T) R (\mu ; T) = R (\lambda ; T) \{(\mu - \lambda) I + (\lambda I - T) \} R (\mu ; T) \\
&= (\mu - \lambda) R (\lambda ; T) R (\mu ; T) + R (\mu ; T).
\end{aligned}
\]
\end{proof}

\begin{theorem}\label{thm:VIII.2.Theorem.3} If $T$ is a bounded linear operator on a complex $B$-space $X$ into $X$, then the following limit exists: 
\begin{equation}\label{eq:VIII.2.3}
\lim\limits _{n \rightarrow \infty} \| T ^{n} \| ^{1/n} = r _{\sigma} (T).
\end{equation}
\end{theorem}

It is called the spectral radius of $T$, and we have 
\begin{equation}\label{eq:VIII.2.4}
\boldsymbol {r} _{\sigma} (T) \leqq \| T \|.
\end{equation}

If $|\lambda| > r_{\sigma}(T)$, then the resolvent $R(\lambda; T)$ exists and is given by the series 
\begin{equation}\label{eq:VIII.2.5}
R (\lambda ; T) = \sum\limits_{n=1} ^{\infty} \lambda^{-n} T ^{n-1}
\end{equation}
which converges in the norm of operators. 

\begin{proof}
 Set $r = \inf\limits_{n \geqq 1} \|T^n\|^{1/n}$. We have to show that $\varlimsup\limits_{n \to \infty} \|T^n\|^{1/n} \leqq r$. For any $\varepsilon > 0$, choose $m$ such that $\|T^m\|^{1/m} \leqq r + \varepsilon$. For arbitrary $n$, write $n = p m + q$ where $0 \leqq q \leqq (m - 1)$. Then, by $\|AB\| \leqq \|A\| \cdot \|B\|$, we obtain 
\[
\left\| T ^{n} \right\| ^{1/n} \leqq \left\| T ^{m} \right\| ^{p/n} \cdot \left\| T \right\| ^{q/n} \leqq (r + \varepsilon) ^{mp/n} \left\| T \right\| ^{q/n}. \]

Since $pm / n \to 1$ and $q / n \to 0$ as $n \to \infty$, we must have $\varlimsup\limits_{n \to \infty} \|T^n\|^{1/n} \leqq r + \varepsilon$. Since $\varepsilon$ was arbitrary, we have proved $\varlimsup\limits_{n \to \infty} \|T^n\|^{1/n} \leqq r$. 

Since $\|T^n\| \leqq \|T\|^n$, we have $\lim\limits_{n \to \infty} \|T^n\|^{1/n} \leqq \|T\|$. The series \eqref{eq:VIII.2.5} is convergent in the norm of operators when $|\lambda| > r_\sigma(T)$. For, if $|\lambda| \geqq r_\sigma(T) + \varepsilon$ with $\varepsilon > 0$, then, by \eqref{eq:VIII.2.3}, $\|\lambda^{-n}T^n\| \leqq (r_\sigma(T) + \varepsilon)^{-n} \cdot (r_\sigma(T) + 2^{-1}\varepsilon)^n$ for large $n$. Multiplication by $(\lambda I - T)$ on the left or right of this series gives $I$ so that the series actually represents the resolvent $R(\lambda; T)$. 
\end{proof}

\begin{corollaryu} The resolvent set $\varrho(T)$ is not empty when $T$ is a bounded linear operator. 
\end{corollaryu}

\begin{theorem}\label{thm:VIII.2.Theorem.4} For a bounded linear operator $T \in L(X, X)$, we have 
\begin{equation}\label{eq:VIII.2.6}
r _{\sigma} (T) = \sup\limits _{\lambda \in \sigma (T)} | \lambda |.
\end{equation}
\end{theorem}

\begin{proof}
 By Theorem~\ref{thm:VIII.2.Theorem.3}, we know that $r_{\sigma}(T) \geqq \sup\limits_{\lambda \in \sigma(T)} |\lambda|$. Hence we have only to show that $r_{\sigma}(T) \leqq \sup\limits_{\lambda \in \sigma(T)} |\lambda|$. 

By Theorem~\ref{thm:VIII.2.Theorem.1}, $R(\lambda; T)$ is holomorphic in $\lambda$ when $|\lambda| > \sup\limits_{\lambda \in \sigma(T)} |\lambda|$. Thus it admits a uniquely determined Laurent expansion in positive and non-positive powers of $\lambda$ convergent in the operator norm for $|\lambda| > \sup |\lambda|$. By Theorem~\ref{thm:VIII.2.Theorem.3}, this Laurent series must coincide with $\sum\limits_{n=1}^{\infty} \lambda^{-n} T^{n-1}$. Hence $\lim\limits_{n \to \infty} \|\lambda^{-n} T^n\| = 0$ if $|\lambda| > \sup\limits_{\lambda \in \sigma(T)} |\lambda|$, and so, for any $\varepsilon > 0$, we must have $\|T^n\| \leqq (\varepsilon + \sup\limits_{\lambda \in \sigma(T)} |\lambda|)^n$ for large $n$. This proves that 
\[
r _{\sigma} (T) = \lim\limits _{n \rightarrow \infty} \| T ^{n} \| ^{1/n} \leqq \sup\limits _{\lambda \in \sigma (T)} | \lambda |.
\]
\end{proof}

\begin{corollaryu} The series $\sum\limits_{n=1}^{\infty} \lambda^{-n} T^{n-1}$ diverges if $|\lambda| < r_{\sigma}(T)$. 
\end{corollaryu}

\begin{proof}
 Let $r$ be the smallest number $\geqq 0$ such that the series $\sum\limits_{n=1}^{\infty} \lambda^{-n} T^{n-1}$ converges in the operator norm for $|\lambda| > r$. The existence of such an $r$ is proved as for ordinary power series in $\lambda^{-1}$. Then, for $|\lambda| > r$, $\lim\limits_{n \to \infty} \|\lambda^{-n} T^n\| = 0$ and so, as in the proof of $r_\sigma(T) \leqq \sup\limits_{\lambda \in \sigma(T)} |\lambda|$, we must have $\lim\limits_{n \to \infty} \|T^n\|^{1/n} \leqq r$. This proves that $r_\sigma(T) \leqq r$. 
\end{proof}


\yossection{3}{The Mean Ergodic Theorem}
\label{sec:VIII.3}

For a particular class of continuous linear operators, the mean ergodic theorem gives a method for obtaining the eigenspace corresponding to the eigenvalue $1$. In this section, we shall state and prove the mean ergodic theorem from the view point of the spectral theory, as was formulated previously by the present author. The historical sketch of the ergodic theory in connection with statistical mechanics will be given in Chapter~\ref{chap:XIII}. 

\begin{theorem}\label{thm:VIII.3.Theorem.1} Let $X$ be a locally convex linear topological space, and $T$ a continuous linear operator on $X$ into $X$. We assume that 
\begin{equation}\label{eq:VIII.3.1}
\parbox{0.92\linewidth}{\everymath{}\let\limits\nolimits the family of operators $\{T^{n}; n=1,2,\ldots\}$ is equicontinuous in the sense that, for any continuous semi-norm $q$ on $X$, there exists a continuous semi-norm $q'$ on $X$ such that $\sup\limits_{n\geqq1}q(T^{n}x)\leqq q'(x)$ for all $x\in X$.}
\end{equation}
Then the closure $R(I - T)^a$ of the range $R(I - T)$ satisfies 
\begin{equation}\label{eq:VIII.3.2}
R (I - T) ^{a} = \left\{x \in X; \lim\limits _{n \rightarrow \infty} T _{n} x = 0, T _{n} = n ^{-1} \sum\limits_{m=1} ^{n} T ^{m} \right\},
\end{equation}
and so, in particular, 
\begin{equation}\label{eq:VIII.3.3}
R (I - T) ^{a} \cap N (I - T) = \{0 \}.
\end{equation}
\end{theorem}

\begin{proof}
 We have $T_{n}(I - T) = n^{-1}(T - T^{n+1})$. Hence, by \eqref{eq:VIII.3.1}, $w \in R(I - T)$ implies that $\lim\limits_{n\to \infty}T_nw = 0$. Next let $z \in R(I - T)^a$. Then, for any continuous semi-norm $q'$ on $X$ and $\varepsilon > 0$, there exists a $w \in R(I - T)$ such that $q'(z - w) < \varepsilon$. Thus, by \eqref{eq:VIII.3.1}, we have $q(T_{n}(z - w)) \leqq n^{-1}\sum\limits_{m=1}^{n}q(T^{m}(z - w)) \leqq q'(z - w) < \varepsilon$. Hence $q(T_{n}z) \leqq q(T_{n}w) + q(T_{n}(z - w)) \leqq q(T_{n}w) + \varepsilon$, and so $\lim\limits_{n\to \infty}T_nz = 0$. This proves that $R(I - T)^a \subseteq \{x \in X; \lim\limits_{n\to \infty}T_nx = 0\}$. 

Let, conversely, $\lim\limits_{n\to \infty}T_nx = 0$. Then, for any continuous semi-norm $q$ on $X$ and $\varepsilon >0$, there exists an $n$ such that $q(x - (x - T_nx)) = q(T_nx) < \varepsilon$. But, by 
\[
\begin{aligned}
&x - T _{n} x  = n ^{-1} \sum\limits_{m=1} ^{n} (I - T ^{m}) x = n ^{-1} \sum\limits_{m=1} ^{n} (I - T) (I + T + T ^{2} + \dots + T ^{m-1}) x,
\end{aligned}
\]
$(x - T_{n}x)\in R(I - T)$. Hence $x$ must belong to $R(I - T)^{a}$ 
\end{proof}

\begin{theorem}[the mean ergodic theorem]\label{thm:VIII.3.Theorem.2} Let condition \eqref{eq:VIII.3.1} be satisfied. Let, for a given $x \in X$, there exist a subsequence $\{n'\}$ of $\{n\}$ 

such that 
\begin{equation}\label{eq:VIII.3.4}
\text{ weak-lim } _ {n ^{\prime} \to \infty} T _ {n ^{\prime}} x = x _{0} \text{ exists. }
\end{equation}
\end{theorem}

Then $Tx_0 = x_0$ and $\lim\limits_{n \to \infty} T_n x = x_0$. 

\begin{proof}
 We have $TT_{n} - T_{n} = n^{-1}(T^{n+1} - T)$, and so, by \eqref{eq:VIII.3.1}, $\lim\limits_{n\to \infty}(TT_nx - T_nx) = 0$. Thus, for any $f\in X'$, $\lim\limits_{n\to \infty}\langle TT_{n'}x,f\rangle = \lim\limits_{n\to \infty}\langle T_{n'}x,T'f\rangle$ exists and $= \lim\limits_{n\to \infty}\langle T_{n'}x,f\rangle = \langle x_0,f\rangle$. Therefore $\langle x_0,f\rangle = \langle Tx_0,f\rangle$ and so, by the arbitrariness of $f\in X'$, we must have $Tx_0 = x_0$. 

We have thus $T^{m}x = T^{m}x_{0} + T^{m}(x - x_{0}) = x_{0} + T^{m}(x - x_{0})$ and so $T_{n}x = x_{0} + T_{n}(x - x_{0})$. But, $(x - x_{0}) = \text{ weak-lim }_{n\to \infty}(x - T_{n'}x)$ and, as proved above, $(x - T_{n'}x)\in R(I - T)$. Therefore, by Theorem~\ref{thm:V.1.Theorem.11} in Chapter~\ref{chap:V}, \ref{sec:V.1}, $(x - x_0)\in R(I - T)^a$. Thus, by Theorem~\ref{thm:VIII.3.Theorem.1}, $\lim\limits_{n\to \infty}T_n(x - x_0) = 0$ and so we have proved that $\lim\limits_{n\to \infty}T_nx = x_0$. 
\end{proof}

\begin{corollaryu} Let condition \eqref{eq:VIII.3.1} be satisfied, and $X$ be locally sequentially weakly compact. Then, for any $x \in X$, $\lim\limits_{n \to \infty} T_n x = x_0$ exists, and the operator $T_0$ defined by $T_0 x = x_0$ is a continuous linear operator such that 
\begin{align}
T _{0}& = T _{0} ^{2} = T T _{0} = T _{0} T, \label{eq:VIII.3.5} \\
R \left(T _{0}\right)& = N (I - T), \label{eq:VIII.3.6} \\
N (T _{0})& = R (I - T) ^{a} = R (I - T _{0}). \label{eq:VIII.3.7}
\end{align}
\end{corollaryu}

Moreover, we have the direct sum decomposition 
\begin{equation}\label{eq:VIII.3.8}
X = R (I - T) ^{a} \oplus N (I - T),
\end{equation}
i.e., any $x \in X$ is represented uniquely as the sum of an element $\in R(I - T)^a$ and an element $\in N(I - T)$. 

\begin{proof}
 The linearity of $T_0$ is clear. The continuity of $T_0$ is proved by the equi-continuity of $\{T_n\}$ implied by \eqref{eq:VIII.3.1}. Next, since $Tx_0 = x_0$, we have $TT_0 = T_0$ and so $T^n T_0 = T_0$, $T_n T_0 = T_0$ which implies that $T_0^2 = T_0$. On the other hand, $T_n - T_n T = n^{-1}(T - T^{n+1})$ and \eqref{eq:VIII.3.1} imply that $T_0 = T_0 T$. The equality \eqref{eq:VIII.3.6} is proved as follows. Let $Tx = x$, then $T^n x = x$, $T_n x = x$. and so $T_0 x = x$, that is, $x \in R(T_0)$. Let conversely, $x \in R(T_0)$. Then, by $T_0^2 = T_0$, we have $T_0 x = x$ and so, by $TT_0 = T_0$, $Tx = TT_0 x = T_0 x = x$. Therefore, the eigenspace of $T$ corresponding to the eigenvalue $1$ of $T$ is precisely the range $R(T_0)$. Hence \eqref{eq:VIII.3.6} is proved. Moreover, we have, by Theorem~\ref{thm:VIII.3.Theorem.1}, $N(T_0) = R(I - T)^a$. But, by $T_0^2 = T_0$, we have $R(I - T_0) \subseteq N(T_0)$, and if $x \in N(T_0)$, then $x = x - T_0 x \in R(I - T_0)$. Thus $N(T_0) = R(I - T_0)$. 

Therefore, by $I = (I - T_0) + T_0$ and \eqref{eq:VIII.3.6} and \eqref{eq:VIII.3.7}, we obtain \eqref{eq:VIII.3.8}. 
\end{proof}

\begin{remarku} The eigenspace $N(\lambda I - T)$ of $T$ belonging to the eigenvalue $\lambda$ with $|\lambda| = 1$ may be obtained as $R(T(\lambda))$, where $T(\lambda)x = \lim\limits_{n\to \infty}n^{-1}\sum\limits_{m=1}^{n}(T / \lambda)^m x$. 
\end{remarku}

The Mean Ergodic Theorem of J. von Neumann. Let $(S, \mathfrak{B}, m)$ be a measure space, and $P$ an equi-measure transformation of $S$, that is, $P$ is a one-one mapping of $S$ onto $S$ such that $P \cdot B \in B$ iff $B$ is $\in B$ and $m(P \cdot B) = m(B)$. Consider the linear operator $T$ on $L^{2}(S, \mathfrak{B}, m)$ onto itself defined by 
\begin{equation}\label{eq:VIII.3.9}
(T x) (s) = x (P s), x \in L ^{2} (S, \mathfrak {B}, m).
\end{equation}

By the equi-measurable property of $P$, we easily see that the operator $T$ is unitary and so the equi-continuity condition \eqref{eq:VIII.3.1} is surely satisfied by $\left\|T^{n}\right\|=1(n=1,2,\ldots)$. Therefore, by the sequential weak compactness of the Hilbert space $L^{2}(S,\mathfrak{B},m)$, we obtain the mean ergodic theorem of J. von Neumann: 
\begin{equation}\label{eq:VIII.3.10}
\begin{aligned}
&\text{ For any } x \in L ^{2} (S, \mathfrak {B}, m), \slim _{n \to \infty} n ^{-1} \sum\limits_{m=1} ^{n} T ^{m} x = x _{0} \in L ^{2} (S, \mathfrak {B}, m) \text{ exists and } T x _{0} = x _{0}.
\end{aligned}
\end{equation}

\begin{remarku} Theorem~\ref{thm:VIII.3.Theorem.1} and Theorem~\ref{thm:VIII.3.Theorem.2} are adapted from K. \textsc{Yosida} \cite{yosida346}. Cf. also S. \textsc{Kakutani} \cite{kakutani125} and F. \textsc{Riesz} \cite{rickart275}. Neumann's mean ergodic theorem was published in J. \textsc{von} \textsc{Neumann} \cite{neumann230}. 
\end{remarku}


\yossection{4}{Ergodic Theorems of the Hille Type Concerning Pseudo-resolvents}
\label{sec:VIII.4}

The notion of resolvent is generalized to that of pseudo-resolvent by E. \textsc{Hille}. We can prove ergodic theorems for pseudo-resolvents by a similar idea to that used in the proof of the mean ergodic theorems in the preceding section. See K. \textsc{Yosida} \cite{yosida347}. Cf. T. \textsc{Kato} \cite{kato134}. These ergodic theorems may be considered as extensions of the abelian ergodic theorems of E. \textsc{Hille} given in E. \textsc{Hille}-R. S. \textsc{Phillips} \cite{hille97}, p. $502$. 

We shall begin with the definition of the pseudo-resolvent. 

\begin{definitionu} Let $X$ be a locally convex complex linear topological space, and $L(X, X)$ the algebra of all continuous linear operators defined on $X$ into $X$. A pseudo-resolvent $J_{\lambda}$ is a function defined on a subset $D(J)$ of the complex $\lambda$-plane with values in $L(X, X)$ such that 
\begin{equation}\label{eq:VIII.4.1}
J _{\lambda} - J _{\mu} = (\mu - \lambda) J _{\lambda} J _{\mu} (\text{ the resolvent equation }).
\end{equation}
\end{definitionu}

\begin{propositionu} All $J_{\lambda}, \lambda \in D(J)$, have a common null space, which we denote by $N(J)$, and a common range which we denote by $R(J)$. Similarly, all $(I - \lambda J_{\lambda}), \lambda \in D(J)$, have a common null space, which we denote by $N(I - J)$, and a common range which we denote by $R(I - J)$. Moreover, we have the commutativity: 
\begin{equation}\label{eq:VIII.4.2}
J _{\lambda} J _{\mu} = J _{\mu} J _{\lambda} (\lambda , \mu \in D (J)).
\end{equation}
\end{propositionu}

\begin{proof}
 By interchanging $\lambda$ and $\mu$ in \eqref{eq:VIII.4.1}, we obtain 
\[
J _{\mu} - J _{\lambda} = (\lambda - \mu) J _{\mu} J _{\lambda} = - (\mu - \lambda) J _{\mu} J _{\lambda},
\]
and hence \eqref{eq:VIII.4.2} is true. The first part of the Proposition is clear from \eqref{eq:VIII.4.1} and \eqref{eq:VIII.4.2}. The second part is also clear from 
\begin{equation}\tag{1'}\label{eq:VIII.4.1'}
(I - \lambda J _{\lambda}) = (I - (\lambda - \mu) J _{\lambda}) (I - \mu J _{\mu})
\end{equation}
which is a variant of \eqref{eq:VIII.4.1}. 
\end{proof}

\begin{theorem}\label{thm:VIII.4.Theorem.1} $A$ pseudo-resolvent $J_{\lambda}$ is a resolvent of a linear operator $A$ iff $N(J) = \{0\}$; and then $R(J)$ coincides with the domain $D(A)$ of $A$. 
\end{theorem}

\begin{proof}
 The ``only if'' part is clear. Suppose $N(J) = \{0\}$. Then, for any $\lambda \in D(J)$, the inverse $J_{\lambda}^{-1}$ exists. We have 
\begin{equation}\label{eq:VIII.4.3}
\lambda I - J _{\lambda} ^{-1} = \mu I - J _{\mu} ^{-1} (\lambda , \mu \in D (J)). \end{equation}

For, by \eqref{eq:VIII.4.1} and \eqref{eq:VIII.4.2}, 
\[
\begin{aligned} J _{\lambda} J _{\mu} (\lambda I - J _{\lambda} ^{-1} - \mu I + J _{\mu} ^{-1}) & = (\lambda - \mu) J _{\lambda} J _{\mu} - J _{\lambda} J _{\mu} (J _{\lambda} ^{-1} - J _{\mu} ^{-1}) \\ & = (\lambda - \mu) J _{\lambda} J _{\mu} - (J _{\mu} - J _{\lambda}) = 0. \end{aligned}
\]

We thus put 
\begin{equation}\label{eq:VIII.4.4}
A = (\lambda I - J _{\lambda} ^{-1}).
\end{equation}

Then $J_{\lambda} = (\lambda I - A)^{-1}$ for $\lambda \in D(J)$. 
\end{proof}
\begin{lemma}\label{lem:VIII.4.Lemma.1} We assume that there exists a sequence $\{\lambda_n\}$ of numbers $\in D(J)$ such that 
\begin{equation}
\label{eq:VIII.4.5}
\begin{minipage}{0.92\linewidth}
\everymath{}\let\limits\nolimits
$\lim\limits_{n\to \infty}\lambda_n = 0$ and the family of operators $\{\lambda_nJ_{\lambda_n}\}$ is equi-continuous.
\end{minipage}
\end{equation}
\end{lemma}


Then we have 
\begin{equation}\label{eq:VIII.4.6}
R (I - J) ^{a} = \left\{x \in X; \lim\limits _{n \rightarrow \infty} \lambda_{n} J _ {\lambda_{n}} x = 0 \right\},
\end{equation}
and hence 
\begin{equation}\label{eq:VIII.4.7}
N (I - J) \cap R (I - J) ^{a} = \{0 \}.
\end{equation}

\begin{proof}
 We have, by \eqref{eq:VIII.4.1}, 
\begin{equation}\label{eq:VIII.4.8}
\lambda J _{\lambda} (I - \mu J _{\mu}) = (1 - \mu (\mu - \lambda) ^{-1}) \lambda J _{\lambda} - \lambda (\lambda - \mu) ^{-1} \mu J _{\mu}. \end{equation}

Hence, by ($5$), the condition $x \in R(I - \mu J_{\mu}) = R(I - J)$ implies that $\lim\limits_{n \to \infty} \lambda_n J_{\lambda_n} x = 0$. Let $y \in R(I - J)^a$. Then, for any continuous semi-norm $q$ on $X$ and $\varepsilon > 0$, there exists an $x \in R(I - J)$ such that $q(y - x) < \varepsilon$. By ($5$), we have, for any continuous semi-norm $q'$ on $X$, 
\[
q ^{\prime} \left(\lambda_{n} J _ {\lambda_{n}} (y - x)\right) \leqq q (y - x) (n = 1, 2, \dots)
\]
with a suitable continuous semi-norm $q$ on $X$. Therefore, by $\lambda_{n}J_{\lambda_{n}}y = \lambda_{n}J_{\lambda_{n}}x + \lambda_{n}J_{\lambda_{n}}(y - x)$, we must have $\lim\limits_{n\to \infty}\lambda_nJ_{\lambda_n}y = 0$. 

Let, conversely, $\lim\limits_{n\to \infty}\lambda_nJ_{\lambda_n}x = 0$. Then, for any continuous semi-norm $q$ on $X$ and $\varepsilon >0$, there exists a $\lambda_{n}$ such that $q(x - (x - \lambda_nJ_{\lambda_n}x)) < \varepsilon$. Hence $x$ must belong to $R(I - \lambda_nJ_{\lambda_n})^a = R(I - J)^a$. 
\end{proof}

\begin{lemman}[Lemma~\ref{lem:VIII.4.Lemma.1}']\yosnamedtag{1'}\label{lem:VIII.4.Lemma.1prime} We assume that there exists a sequence $\{\lambda_n\}$ of numbers $\in D(J)$ such that 
$\lim\limits_{n\to\infty}\left|\lambda_{n}\right|=\infty$ and the family of operators $\{\lambda_{n}J_{\lambda_{n}}\}$ is equi-continuous. ($5$') Then we have 
\begin{equation}\tag{6'}\label{eq:VIII.4.6'}
R (J) ^{a} = \left\{x \in X; \lim\limits _{n \rightarrow \infty} \lambda_{n} J _ {\lambda_{n}} x = x \right\},
\end{equation}
and hence 
\begin{equation}\tag{7'}\label{eq:VIII.4.7'}
N (J) \cap R (J) ^{a} = \{0 \}.
\end{equation}
\end{lemman}

\begin{proof}
 We have, by \eqref{eq:VIII.4.1}, 
\[
\lambda J _{\lambda} J _{\mu} = \frac {\mu}{\lambda} \lambda J _{\lambda} J _{\mu} - \frac {1}{\lambda} \lambda J _{\lambda} + J _{\mu}. \]

Hence, by ($5$'), the condition $x \in R(J_{\mu}) = R(J)$ implies that $\lim\limits_{n \to \infty} \lambda_n J_{\lambda_n} x = x$. Let $y \in R(J)^a$. Then, for any continuous semi-norm $q$ on $X$ and $\varepsilon > 0$, there exists an $x \in R(J)$ such that $q(y - x) < \varepsilon$. By ($5$'), we have, for any continuous semi-norm $q'$ on $X$, 
\[
q ^{\prime} \left(\lambda_{n} J _ {\lambda_{n}} (y - x)\right) \leqq q (y - x) \quad (n = 1, 2, \dots)
\]
with a suitable continuous semi-norm $q$ on $X$. Therefore, by $(5')$ and 
\[
\lambda_{n} J _ {\lambda_{n}} y - y = \left(\lambda_{n} J _ {\lambda_{n}} x - x\right) + (x - y) + \lambda_{n} J _ {\lambda_{n}} (y - x),
\]
we must have $\lim\limits_{n\to \infty}\lambda_nJ_{\lambda_n}y = y.$ 

Let, conversely, $\lim\limits_{n\to \infty}\lambda_nJ_{\lambda_n}x = x$. Then, for any continuous semi-norm $q$ on $X$ and $\varepsilon >0$, there exists a $\lambda_{n}$ such that $q(x - \lambda_nJ_{\lambda_n}x) < \varepsilon$. Hence $x$ must belong to $R(J_{\lambda_n})^a = R(J)^a$. 
\end{proof}

\begin{theorem}\label{thm:VIII.4.Theorem.2} Let ($5$) be satisfied. Let, for a given $x \in X$, there exist a subsequence $\{n'\}$ of $\{n\}$ such that 
\begin{equation}\label{eq:VIII.4.9}
\underset {n \to \infty} {\text{ weak-lim }} \lambda_ {n ^{\prime}} J _ {\lambda_ {n ^{\prime}}} x = x _{h} \text{ exists. }
\end{equation}
\end{theorem}

Then $x_{h} = \lim\limits_{n\to \infty}\lambda_{n}J_{\lambda_{n}}x$ and $x_{h}\in N(I - J), x_{p} = (x - x_{h})\in R(I - J)^{a}$. 

\begin{proof}
 Setting $\mu = \lambda_{n'}$ in \eqref{eq:VIII.4.1} and letting $n' \to \infty$, we see, by ($5$), that $(I - \lambda J_{\lambda}) x = (I - \lambda J_{\lambda})(x - x_h)$, that is, $(I - \lambda J_{\lambda}) x_h = 0$. Hence $x_h \in N(I - J)$ and so 
\begin{equation}\label{eq:VIII.4.10}
\lambda_{n} J _ {\lambda_{n}} x = x _{h} + \lambda_{n} J _ {\lambda_{n}} (x - x _{h}). \end{equation}

Therefore we have only to prove that $\lim\limits_{n\to \infty}\lambda_nJ_{\lambda_n}(x - x_h) = 0$, or, by Lemma~\ref{lem:VIII.4.Lemma.1}, $(x - x_h)\in R(I - J)^a$. But $(x - \lambda_nJ_{\lambda_n}x)\in R(I - J)$, and so, by Theorem~\ref{thm:V.1.Theorem.11} in Chapter~\ref{chap:V}, \ref{sec:V.1}, we must have $(x - x_h)\in R(I - J)^a$. 
\end{proof}

\begin{corollary}\label{cor:VIII.4.Corollary.1} Let ($5$) be satisfied, and let $X$ be locally sequentially weakly compact. Then 
\begin{equation}\label{eq:VIII.4.11}
X = N (I - J) \oplus R (I - J) ^{a}.
\end{equation}
\end{corollary}

\begin{proof}
 For any $x \in X$, let $x_{h} = \lim\limits_{n \to \infty} \lambda_{n} J_{\lambda_{n}} x$ and $x_{p} = (x - x_{h})$ be the components of $x$ in $N(I - J)$ and $R(I - J)^{a}$, respectively. 
\end{proof}

\begin{theoremn}[Theorem~\ref{thm:VIII.4.Theorem.2}']\yosnamedtag{2'}\label{thm:VIII.4.Theorem.2prime} Let $(5')$ be satisfied. Let, for a given $x \in X$, there exist a subsequence $\{n'\}$ of $\{n\}$ such that 
\begin{equation}\tag{9'}\label{eq:VIII.4.9'}
\underset {n \rightarrow \infty} {\text{ weak-lim }} \lambda_ {n ^{\prime}} J _ {\lambda_ {n ^{\prime}}} x = x _ {h ^{\prime}} \text{ exists. }
\end{equation}

Then $x_{h'} = \lim\limits_{n\to \infty}\lambda_nJ_{\lambda_n}x$ and $x_{h'}\in R(J)^a, x_{p'} = (x - x_{h'})\in N(J).$
\end{theoremn}

\begin{proof}
 Setting $\mu = \lambda_{n'}$ in \eqref{eq:VIII.4.8} and letting $n' \to \infty$, we see, by ($5$'), that $\lambda J_{\lambda}(x - x_{h'}) = 0$, that is, $(x - x_{h'}) \in N(J)$. Hence 
\begin{equation}\tag{10'}\label{eq:VIII.4.10'}
\lambda_{n} J _ {\lambda_{n}} x = \lambda_{n} J _ {\lambda_{n}} x _ {h ^{\prime}}. \end{equation}

Therefore we have only to prove that $\lim\limits_{n\to \infty}\lambda_nJ_{\lambda_n}x_{h'} = x_{h'}$, or, by Lemma~\ref{lem:VIII.4.Lemma.1prime}, $x_{h'}\in R(J)^a$. But $\lambda_{n'}J_{\lambda_n},x_{h'}\in R(J)$, and so, by Theorem~\ref{thm:V.1.Theorem.11} in Chapter~\ref{chap:V}, \ref{sec:V.1}, we must have $x_{h'}\in R(J)^a$. 
\end{proof}

\begin{corollaryn}[Corollary~\ref{cor:VIII.4.Corollary.1}']\yosnamedtag{1'}\label{cor:VIII.4.Corollary.1prime} Let $(5')$ be satisfied, and let $X$ be locally sequentially weakly compact. Then 
\begin{equation}\tag{11'}\label{eq:VIII.4.11'}
X = N (J) \oplus R (J) ^{a}.
\end{equation}
\end{corollaryn}

\begin{proof}
 For any $x \in X$, let $x_{h'} = \lim\limits_{n \to \infty} \lambda_n J_{\lambda_n} x$ and $x_{p'} = (x - x_{h'})$ be the components of $X$ in $R(J)^a$ and $N(J)$, respectively. 
\end{proof}

\begin{remarku} As a Corollary we obtain: In a reflexive $B$-space $X$, a pseudo-resolvent $J_{\lambda}$ satisfying $(5')$ is the resolvent of a closed linear operator $A$ with dense domain iff $R(J)^{a} = X$. This result is due to $T$. \textsc{Kato}, loc. cit. The proof is easy, since, by Eberlein's theorem, a $B$-space $X$ is locally sequentially weakly compact iff $X$ is reflexive. 
\end{remarku}


\yossection{5}{The Mean Value of an Almost Periodic Function}
\label{sec:VIII.5}

As an application of the mean ergodic theorem we shall give an existence proof of the mean value of an almost periodic function. 

\begin{definition}\label{def:VIII.5.Definition.1} $A$ set $G$ of elements $g, h, \ldots$ is called a group if in $G$ a product (in general non-commutative) $gh$ of any pair $(g, h)$ of elements $\in G$ is defined satisfying the following conditions: 
\begin{align}
g h& \in G, \label{eq:VIII.5.1} \\
(g h) k& = g (h k) \text{ (the associativity) }, \label{eq:VIII.5.2}
\end{align}
\begin{equation}\label{eq:VIII.5.3}
\begin{minipage}[b]{0.92\linewidth}
there exists a unique element $e$ in $G$ such that $eg = ge = g$ for all $g \in G$; $e$ is called the identity element of the group $G$,
\end{minipage}
\end{equation}

\begin{equation}\label{eq:VIII.5.4}
\begin{minipage}[b]{0.92\linewidth}
for every element $g \in G$, there exists a uniquely determined element in $G$, which is denoted by $g^{-1}$, such that $gg^{-1} = g^{-1}g = e$; the element $g^{-1}$ is called the inverse element of $g$.
\end{minipage}
\end{equation}
\end{definition}


Clearly, $g$ is the inverse of $g^{-1}$ so that $(g^{-1})^{-1} = g$. A group $G$ is said to be commutative if $gh = hg$ for all $g, h \in G$. 

\begin{examplen} The totality of complex matrices of order $n$ with determinants equal to unity is a group with respect to the matrix multiplication; it is called the complex unimodular group of order $n$. The identity of this group is the identity matrix and the inverse element of the matrix $a$ is the inverse matrix $a^{-1}$. The real unimodular group is defined analogously. These groups are non-commutative when $n \geqq 2$.

\end{examplen}
\begin{definition}[J. \textsc{von} \textsc{Neumann} {\cite{neumann231}}]\label{def:VIII.5.Definition.2} Given an abstract group $G$. A complex-valued function $f(g)$ defined on $G$ is called almost periodic on $G$ if the following condition is satisfied: 

\begin{equation}\label{eq:VIII.5.5}
\begin{minipage}[b]{0.92\linewidth}
the set of functions $\{g_s(f, h); s \in G\}$, where $f_s(g, h) = f(gh)$, defined on the direct product $G \times G$ is totally bounded with respect to the topology of uniform convergence on $G \times G$.
\end{minipage}
\end{equation}
\end{definition}


\begin{examplen} Let $G$ be the set $\mathbb{R}^{1}$ of all real numbers in which the group multiplication is defined as the addition of real numbers; this group $\mathbb{R}^{1}$ is called the additive group of real numbers. The function $f(g) = e^{i\alpha g}$, where $\alpha$ is a real number and $i = \sqrt{-1}$, is almost periodic on $\mathbb{R}^{1}$. This we easily see from the addition theorem $f(gsh) = e^{i\alpha g} e^{i\alpha s} e^{i\alpha h}$ and the fact that $\{e^{i\alpha t}; t \in \mathbb{R}^{1}\}$ is totally bounded as a set of complex numbers of absolute value $1$.

\end{examplen}
\begin{proposition}\label{prop:VIII.5.Proposition.1} Suppose $f(g)$ is an almost periodic function on $G$. If we define, following A. \textsc{Weil}, 
\begin{equation}\label{eq:VIII.5.6}
\dist (s, u) = \sup\limits _{g, h \in G} | f (g s h) - f (g u h) |,
\end{equation}
then 
\begin{equation}\label{eq:VIII.5.7}
\dist (s, u) = \dist (a s b, a u b).
\end{equation}
\end{proposition}

\begin{proof}
 Clear from the definition of the group. 
\end{proof}

\begin{corollary}\label{cor:VIII.5.Corollary.1} The set $E$ of all elements $s$ which satisfies $\mathrm{dis}(s, e) = 0$ constitutes an invariant subgroup in $G$, that is, we have 

\begin{equation}\label{eq:VIII.5.8}
\begin{minipage}[b]{0.92\linewidth}
if $e_1, e_2 \in E$, then $e_1e_2 \in E$ and $ae_1a^{-1} \in E$ for every $a \in G$.
\end{minipage}
\end{equation}
\end{corollary}


\begin{proof}
 Let $\dist(e_1, e) = 0$, $\dist(e_2, e) = 0$. Then, by \eqref{eq:VIII.5.7} and the triangle inequality, we obtain 
\[
\dist \left(e _{1} e _{2}, e\right) \leqq \dist \left(e _{1} e _{2}, e _{1} e\right) + \dist \left(e _{1} e, e\right) = 0 + 0 = 0. \]

Similarly we have $\mathrm{dis}(ae_1a^{-1},e) = \mathrm{dis}(ae_1a^{-1},aea^{-1}) = 0$ from $\mathrm{dis}(e_1,e) = 0$. 
\end{proof}

\begin{corollary}\label{cor:VIII.5.Corollary.2} If we write $s \equiv u (\bmod E)$ when $su^{-1} \in E$, then $s \equiv u (\bmod E)$ is equivalent to $\dist(s, u) = 0$. 
\end{corollary}

\begin{proof}
 Clear from $\dist(su^{-1}, e) = \dist(s, eu) = \dist(s, u)$. 
\end{proof}

\begin{corollary}\label{cor:VIII.5.Corollary.3} The concept $s \equiv u \pmod{E}$ has all the general properties of equivalence, namely 
\begin{align}
s& \equiv s (\mathrm{mod} E), \label{eq:VIII.5.9} \\
s& \equiv u (\mathrm{mod} E), \text{ then } u \equiv s (\mathrm{mod} E), \label{eq:VIII.5.10} \\
& \text{ if } s _{1}& \equiv s _{2} (\bmod E) \text{ and } s _{2} \equiv s _{3} (\bmod E), \text{ then } s _{1} \equiv s _{3} \notag
\\
& (\bmod E). \label{eq:VIII.5.11}
\end{align}
\end{corollary}

\begin{proof}
 Clear from Corollary~\ref{cor:VIII.5.Corollary.2} and the triangle inequality for the $\dist(s, u)$. 
\end{proof}

Hence, as in the case of the factor space in a linear space, we can define the factor group or residue class group $G$/$E$ as follows: we shall denote the set of all elements $\in G$ equivalent (mod $E$) to a fixed element $x \in G$ by $\xi_{x}$, the residue class (mod $E$) containing $x$; then the set of all residue classes $\xi_{x}$ constitutes a group $G$/$E$ by the notion of the product 
\begin{equation}\label{eq:VIII.5.12}
\xi_{x} \xi_{y} = \xi_{xy}.
\end{equation}

To justify this definition \eqref{eq:VIII.5.12} of the product, we have to show that 

\begin{equation}\label{eq:VIII.5.13}
\begin{minipage}[b]{0.92\linewidth}
if $x_{1} \equiv x_{2} (\bmod E)$, $y_{1} \equiv y_{2} (\bmod E)$, then $x_{1}y_{1} \equiv x_{2}y_{2} (\bmod E)$.
\end{minipage}
\end{equation}

This is clear, since we have by \eqref{eq:VIII.5.7} and Corollary~\ref{cor:VIII.5.Corollary.2}, 
\[
\begin{aligned}
&\mathrm{dis} (x _{1} y _{1}, x _{2} y _{2})  \leqq \mathrm{dis} (x _{1} y _{1}, x _{2} y _{1}) + \mathrm{dis} (x _{2} y _{1}, x _{2} y _{2}) = \mathrm{dis} (x _{1}, x _{2}) + \mathrm{dis} (y _{1}, y _{2}) = 0 + 0 = 0.
\end{aligned}
\]

Since the function $f(x)$ takes the same constant value on the residue class $\xi_{x}$, we may consider $f(x)$ as a function $F(\xi_{x})$ defined on the residue class group $G$/$E$. 

\begin{corollary}\label{cor:VIII.5.Corollary.4} The residue class group is a metric space by the distance 
\begin{equation}\label{eq:VIII.5.14}
\dist \left(\xi_{x}, \xi_{y}\right) = \dist (x, y).
\end{equation}
\end{corollary}

\begin{proof}
 $x \equiv x_{1} (\bmod E)$ and $y \equiv y_{1} (\bmod E)$ imply 
\[
\dist (x, y) \leqq \dist (x, x _{1}) + \dist (x _{1}, y _{1}) + \dist (y _{1}, y) = 0 + \dist (x _{1}, y _{1}) + 0
\]
and $\dist(x_1, y_1) \leqq \dist(x, y)$ to the effect that $\dist(x, y) = \dist(x_1, y_1)$. Thus \eqref{eq:VIII.5.14} defines a distance in $G/E$. 
\end{proof}

\begin{corollary}\label{cor:VIII.5.Corollary.5} The group $G$/$E$ is a topological group with respect to the distance $\text{ dis }(\xi_{x}, \xi_{y})$, that is, the operation of multiplication $\xi_{x}\xi_{y}$ is continuous as a mapping from the product space $(G/E)\times(G/E)$ onto $G$/$E$, and the operation $\xi_{x}^{-1}$ is continuous as a mapping from $G$/$E$ onto $G$/$E$. 
\end{corollary}

\begin{proof}
 We have, by \eqref{eq:VIII.5.7}, 
\begin{gather}\quad \dist (s u, s ^{\prime} u ^{\prime})\leqq \dist (s u, s ^{\prime} u) + \dist (s ^{\prime} u, s ^{\prime} u ^{\prime}) = \dist (s, s ^{\prime}) + \dist (u, u ^{\prime}) \text{ and } \notag
\\\dist \left(s ^{-1}, u ^{-1}\right) = \dist \left(s s ^{-1} u, s u ^{-1} u\right) = \dist (u, s) = \dist (s, u). \notag
\end{gather}
\end{proof}

We have thus proved the following 

\begin{theorem}[A. \textsc{Weil}]\label{thm:VIII.5.Theorem.1} The topological group $G / E$, metrized by \eqref{eq:VIII.5.14}, is totally bounded, and the function $f(x)$ gives rise to a function $F(\xi_x) (= f(x))$ which is uniformly continuous on this group $G / E$. 
\end{theorem}

\begin{proof}
 The uniform continuity of the function $F(\xi_x)$ is clear from 
\[
\left| F \left(\xi_{x}\right) - F \left(\xi_{y}\right) \right| = \left| f (x) - f (y) \right| \leqq \dist (x, y) = \dist \left(\xi_{x}, \xi_{y}\right). \]

The almost periodicity of the function $f(x)$ implies, by \eqref{eq:VIII.5.7} and \eqref{eq:VIII.5.14}, that the metric space $G / E$ is totally bounded. 
\end{proof}

By the above theorem, the theory of almost periodic functions is reduced to the study of a uniformly continuous function $f(g)$, defined on a totally bounded topological group $G$, metrized by a metric $\text{ dis }(g_{1}, g_{2})$ satisfying condition \eqref{eq:VIII.5.7}. By virtue of this fact, we shall give a proof for the existence of the mean value of an almost periodic function. 

Since $G$ is totally bounded, there exists, for any $\varepsilon > 0$, a finite system of points $g_1, g_2, \ldots, g_n$ such that $\min\limits_{1 \leqq i \leqq n} \text{ dis } (g, g_i) \leqq \varepsilon$ for any $g \in G$. Hence, collecting these finite systems of points corresponding to $\varepsilon = 1, 2^{-1}, 3^{-1}, \ldots$, respectively, we see that there exists a countable system $\{g_j\}$ of points $\in G$ such that $\{g_j\}$ is dense in $G$. We take a sequence of positive numbers $\alpha_j$ such that $\sum\limits_{j=1}^{\infty} \alpha_j = 1$. Let $C(G)$ be the set of all uniformly continuous complex-valued functions $h(g)$ defined on $G$. $C(G)$ is a $B$-space by the operation of the function sum and the norm $\|h\| = \sup |h(g)|$. We define an operator $T$ defined on $C(G)$ into $C(G)$ by $g \in G$. 
\begin{equation}\label{eq:VIII.5.15}
(T \cdot h) (g) = \sum\limits_{j=1} ^{\infty} \alpha_{j} h (g _{j} g).
\end{equation}

By the uniform continuity of $h(g)$ on $G$, there exists, for any $\varepsilon > 0$, a $\delta > 0$ such that $\mathrm{dis}(g, g') \leqq \delta$ implies $|h(g) - h(g')| \leqq \varepsilon$. Thus, by \eqref{eq:VIII.5.7}, $|h(g_j g) - h(g_j g')| \leqq \varepsilon (j = 1, 2, \ldots)$ whenever $\mathrm{dis}(g, g') \leqq \delta$. Hence it is easy to see, by $\alpha_j > 0$ and $\sum\limits_{j=1}^{\infty} \alpha_j = 1$, that $T$ is a bounded linear operator on $C(G)$ into $C(G)$. By the same reasoning we see that the set of functions $h_n(g)$ defined by 
\begin{align}
h _{n} (g)& = n ^{-1} \sum\limits_{m=1} ^{n} \left(T ^{m} h\right) (g), \text{ which is of the form } \notag \\
h _{n} (g)& = \sum\limits_{j=1} ^{\infty} \beta_{j} h \left(g _{j} ^{\prime} g\right) \quad \text{ with } \quad \beta_{j} > 0, \sum\limits_{j=1} ^{\infty} \beta_{j} = 1, \label{eq:VIII.5.16}
\end{align}
is equi-bounded and equi-continuous with respect to $n$. Hence, by the Ascoli-Arzelà theorem, the sequence $\{h_n(g)\}$ contains a subsequence which is uniformly convergent on $G$. 

Therefore, by the mean ergodic theorem, there exists an $h^*(g) \in C(G)$ such that 
\begin{equation}\label{eq:VIII.5.17}
\lim\limits _{n \rightarrow \infty} \sup\limits _{g} | h _{n} (g) - h ^{*} (g) | = 0 \text{ and } T h ^{*} = h ^{*}.
\end{equation}

\begin{proposition}\label{prop:VIII.5.Proposition.2} $h^{*}(g)$ is identically equal to a constant. 
\end{proposition}

\begin{proof}
 We may assume, without losing the generality, that $h(g)$ and $h^{*}(g)$ are real-valued. Suppose that there exist a point $g_{0} \in G$ and a positive constant $\delta$ such that 
\[
h ^{*} (g _{0}) \leqq \beta - 2 \delta , \text{ where } \beta = \sup\limits _{g \in G} h ^{*} (g). \]

By the continuity of $h^*(g)$, there exists a positive number $\varepsilon$ such that $\text{ dis }(g', g'') \leqq \varepsilon$ implies $|h^*(g') - h^*(g'')| \leqq \delta$; in particular, we have $h^*(g'') \leqq \beta - \delta$ whenever $\text{ dis }(g_0, g'') \leqq \varepsilon$. Since the sequence $\{g_j\}$ is dense in $G$, there exists, for any $\varepsilon > 0$, an index $n$ such that, for any $g \in G$, we have $\min\limits_{1 \leqq i \leqq n} \text{ dis }(g, g_j) \leqq \varepsilon$. Hence, by \eqref{eq:VIII.5.7}, we have, for any $g \in G$, 
\[
\min\limits _{1 \leqq j \leqq n} \dist (g _{0}, g _{j} g) \leqq \varepsilon .
\]

Let the minimum be attained at $j = j_0$. Then 
\[
h ^{*} (g) = \left(T h ^{*}\right) (g) = \sum\limits_{j=1} ^{\infty} \alpha_{j} h ^{*} \left(g _{j} g\right) \leqq \alpha_ {j _{0}} (\beta - \delta) + \left(1 - \alpha_ {j _{0}}\right) \beta = \beta - \alpha_ {j _{0}} \delta < \beta ,
\]
contrary to the assumption that $g$ was an arbitrary point of $G$. Therefore $h^{*}(g)$ is identically equal to a constant. 
\end{proof}

\begin{definition}\label{def:VIII.5.Definition.3} We shall call the constant value $h^{*}(g)$ the left mean value of $h(g)$ and denote its value by $M_{g}^{l}(h(g))$: 
\begin{equation}\label{eq:VIII.5.18}
M _{g} ^{l} (h (g)) = \lim\limits _{n \rightarrow \infty} n ^{-1} \sum\limits_{m=1} ^{n} \left(T ^{m} h\right) (g).
\end{equation}
\end{definition}

\begin{theorem}[J. von Neumann]\label{thm:VIII.5.Theorem.2} We have 
\begin{align}
M _{g} ^{l} (\alpha h (g))& = \alpha M _{g} ^{l} (h (g)), \tag{i}\label{eq:VIII.5.i} \\
M _{g} ^{l} (h _{1} (g) + h _{2} (g))& = M _{g} ^{l} (h _{1} (g)) + M _{g} ^{l} (h _{2} (g)), \tag{ii}\label{eq:VIII.5.ii} \\
M _{g} ^{l} \eqref{eq:VIII.5.1}& = 1, \tag{iii}\label{eq:VIII.5.iii} \\
\text{ if } h (g)& \geqq 0 \text{ on } G, \text{ then } M _{g} ^{l} (h (g)) \geqq 0; \notag \\
\text{ if, moreover, } h (g) \not& \equiv 0, \text{ then } M _{g} ^{l} (h (g)) > 0, \tag{iv}\label{eq:VIII.5.iv} \\
\left| M _{g} ^{l} (h (g)) \right|& \leqq M _{g} ^{l} (\left| h (g) \right|), \tag{v}\label{eq:VIII.5.v} \\
M _{g} ^{l} \overline {{(h (g))}}& = \overline {{M _{g} ^{l} (h (g))}}, \tag{vi}\label{eq:VIII.5.vi} \\
M _{g} ^{l} (h (g a))& = M _{g} ^{l} (h (g)), \tag{vii}\label{eq:VIII.5.vii} \\
M _{g} ^{l} (h (a g))& = M _{g} ^{l} (h (g)), \tag{vii'}\label{eq:VIII.5.vii'} \\
M _{g} ^{l} (h (g ^{-1}))& = M _{g} ^{l} (h (g)). \tag{viii}\label{eq:VIII.5.viii}
\end{align}
\end{theorem}

\begin{proof}
 By definition \eqref{eq:VIII.5.18}, it is clear that (i), (ii), (iii), the first part of (iv), (v) and (vi) are true. The truth of (vii) is proved by Proposition~\ref{prop:VIII.5.Proposition.2}. (vii') is proved as follows.


Starting with the linear operator $T'$ defined by 
\[
\left(T ^{\prime} h\right) (g) = \sum\limits_{j=1} ^{n} \alpha_{j} h \left(g g _{j}\right),
\]
we can also define a right mean $M_g^r(h(g))$ which, as a functional of $h(g)$, satisfies (i), (ii), (iii), the first part of (iv), (v), (vi) and (vii'). We thus have to prove that the left mean $M_g^l(h(g))$ coincides with the right mean $M_g^r(h(g))$. By its definition of the left mean, there exist, for any $\varepsilon > 0$, a sequence of elements $\{k_j\} \subseteq G$ and a sequence of positive numbers $\beta^j$ with $\sum\limits_{j=1}^{\infty} \beta_j = 1$ such that
\begin{equation}\label{eq:VIII.5.19}
\sup\limits _{g} \left| \sum\limits_{j=1} ^{\infty} \beta_{j} h (k _{j} g) - M _{g} ^{l} (h (g)) \right| \leqq \varepsilon .
\end{equation}

Similarly, there exist a sequence of elements $\{s_j\} \subseteq G$ and a sequence of positive numbers $\gamma_j$ with $\sum\limits_{j=1}^{\infty} \gamma_j = 1$ such that 
\begin{equation}\label{eq:VIII.5.20}
\sup\limits _{g} \left| \sum\limits_{j=1} ^{\infty} \gamma_{j} h (g s _{j}) - M _{g} ^{\prime} (h (g)) \right| \leqq \varepsilon .
\end{equation}

We have, from \eqref{eq:VIII.5.19} and (vii) 
\[
\sup\limits _{g} \left| \sum\limits_{i,j} \gamma_{i} \beta_{j} h (k _{j} g s _{i}) - M _{g} ^{l} (h (g)) \right| \leqq \varepsilon ,
\]
and, similarly from \eqref{eq:VIII.5.20} and $(\mathrm{vii}^{\prime})$ 
\[
\sup\limits _{g} \left| \sum\limits_{i,j} \gamma_{i} \beta_{j} h (k _{j} g s _{i}) - M _{g} ^{r} (h (g)) \right| \leqq \varepsilon .
\]

Hence we must have $M_g^l(h(g)) = M_g'(h(g))$. 

We next remark that a linear functional $M_{g}(h(g))$ defined on $C(G)$ is uniquely determined by the properties (i), (ii), (iii), the first part of (iv), (v), (vi) and (vii) (or (vii') as well). In fact, we have, by \eqref{eq:VIII.5.20},
\[
M _{g} ^{\prime} (h (g)) - \varepsilon \leqq \sum\limits_{i=1} ^{\infty} \gamma_{i} h (g s _{i}) \leqq M _{g} ^{\prime} (h (g)) + \varepsilon \quad \text{ for real-valued } h (g)
\]

Hence, for real-valued $h(g)$, $M_g(h(g))$ must coincide with the right mean $M_g^r(h(g))$ and hence with the left mean $M_g^l(h(g))$ as well. Therefore we see that we must have $M_g = M_g' = M_g^l$. Being equal to the right mean, $M_g$ must satisfy (vii'). Moreover, since $M_g^l(h(g^{-1}))$ satisfies, as a linear functional, (i), (ii), (iii), the first part of (iv), (v), (vi) and (vii'), we must have $M_g^l(h(g^{-1})) = M_g^r(h(g)) = M_g^l(h(g))$.

Finally we shall prove the last part of (iv). Suppose $h(g_{0}) > 0$. For any $\varepsilon > 0$, there exists, by the total boundedness of $G$, a finite system of elements $s_{1}, s_{2}, \ldots, s_{n}$ such that 
\[
\min\limits _{1 \leqq i \leqq n} \sup\limits _{g} | h (g s _{i}) - h (g s) | < \varepsilon
\]
for all $s \in G$. This we see by the uniform continuity of $h(g)$ and the fact that $\dist(gs_i, gs) = \dist(s_i, s)$. Hence, for $\varepsilon = h(g_0)/2$, we obtain, for any $s \in G$, a suffix $i, 1 \leqq i \leqq n$, such that 
\[
h \left(g _{0} s _{i} ^{-1} s\right) \geqq h \left(g _{0}\right) / 2.
\]

Thus, by the non-negativity of the function $h(g)$, we obtain 
\[
\sum\limits_{i=1} ^{n} h \left(g _{0} s _{i} ^{-1} s\right) \geqq h \left(g _{0}\right) / 2 > 0 \quad \text{ for all } \quad s \in G.
\]

Therefore, by taking the right mean of both sides, we have 
\[
M _{g} ^{r} \left(\sum\limits_{i=1} ^{n} h \left(g _{0} s _{i} ^{-1} s\right)\right) = n M _{s} ^{r} (h (s)) \geqq h \left(g _{0}\right) / 2 > 0.
\]
\end{proof}

\begin{remarku} The happy idea of introducing the distance \eqref{eq:VIII.5.1} is shown in A. \textsc{Weil} \cite{weil334}. The application of the mean ergodic theorem to the existence proof of the mean value is due to the present author. See also W. \textsc{Maak} \cite{lumer185}. 
\end{remarku}


\yossection{6}{The Resolvent of a Dual Operator}
\label{sec:VIII.6}

\begin{lemma}\label{lem:VIII.6.Lemma.1} Let $X$ and $Y$ be complex $B$-spaces. Let $T$ be a linear operator with $D(T)^a = X$ and $R(T) \subseteq Y$. Then $(T')^{-1}$ exists iff $R(T)^a = Y$.
\end{lemma}

\begin{proof}
If $T' \nu_0^* = 0$, then 
\[
\langle x, T ^{\prime} y _{0} ^{*} \rangle = \langle T x, y _{0} ^{*} \rangle = 0 \quad \text{ for all } \quad x \in D (T),
\]
and hence $y_0^*(R(T)^a) = 0$. Thus $R(T)^a = Y$ implies $y_0^* = 0$ and so $T'$ has an inverse. On the other hand, if $y_0 \bar{\in} R(T)^a$, then, the Hahn-Banach theorem asserts that there exists a continuous linear functional $y_0^* \in Y'$ such that $y_0^*(y_0) = 1$ and $y_0^*(R(T)^a) = 0$. Hence $\langle Tx, y_0^* \rangle = 0$ for all $x \in D(T)$, and so $y_0^* \in D(T')$ and $T'y_0^* = 0$, whereas $y_0^*(y_0) \neq 0$, i.e., $y_0^* \neq 0$. Therefore, the condition $R(T)^a \neq Y$ implies that $T'$ cannot have an inverse. 
\end{proof}

\begin{theorem}[R. S. \textsc{Phillips} {\cite{phillips252}}]\label{thm:VIII.6.Theorem.1} Let $T$ be a linear operator with an inverse and such that $D(T)^{a} = X$ and $R(T)^{a} = Y$, where $X$ and $Y$ are $B$-spaces. Then 
\begin{equation}\label{eq:VIII.6.1}
(T ^{\prime}) ^{-1} = (T ^{-1}) ^{\prime}.
\end{equation}
$T^{-1}$ is bounded on $Y$ iff $T$ is closed and $(T')^{-1}$ is bounded on $X_s'$. 
\end{theorem}

\begin{proof}
 $(T^{-1})'$ exists because $D(T^{-1}) = R(T)$ is dense in $Y$. $(T')^{-1}$ exists by Lemma~\ref{lem:VIII.6.Lemma.1}. We have to show the equality \eqref{eq:VIII.6.1}. If $y \in R(T)$ and $y^* \in D(T')$, then 
\[
\langle y, y ^{*} \rangle = \langle T T ^{-1} y, y ^{*} \rangle = \langle T ^{-1} y, T ^{\prime} y ^{*} \rangle . \]

Hence $R(T') \subseteq D((T^{-1})')$ and $(T^{-1})'(T'y^*) = y^*$ for all $y^* \in D(T')$. Thus $(T^{-1})'$ is an extension of $(T')^{-1}$. Next, if $x \in D(T)$, then 
\[
\langle x, x ^{*} \rangle = \langle T ^{-1} T x, x ^{*} \rangle = \langle T x, (T ^{-1}) ^{\prime} x ^{*} \rangle \quad \text{ for all } \quad x ^{*} \in D ((T ^{-1}) ^{\prime}).
\]

Hence $R((T^{-1})') \subseteq D(T')$ and $T'(T^{-1})' x^* = x^*$ for all $x^* \in D((T^{-1})')$. Thus $(T^{-1})'$ is a contraction of $(T')^{-1}$. Therefore we have proved \eqref{eq:VIII.6.1}. 

If, in addition, $T^{-1}$ is bounded on $Y$, then $(T^{-1})'$ is also bounded. Conversely, if $(T')^{-1}$ is bounded on $X_s'$, then, for all $x \in R(T)$ and $x^* \in X'$, we have, by \eqref{eq:VIII.6.1}, 
\[
\begin{aligned}
&| \langle T ^{-1} x, x ^{*} \rangle |  = | \langle x, (T ^{-1}) ^{\prime} x ^{*} \rangle | = | \langle x, (T ^{\prime}) ^{-1} x ^{*} \rangle | \leqq \| (T ^{\prime}) ^{-1} \| \cdot \| x ^{*} \| \cdot \| x \|.
\end{aligned}
\]

Since $T^{-1}$ is closed and $R(T)^{a} = Y$, $T^{-1}$ must be bounded. 
\end{proof}
\begin{lemma}\label{lem:VIII.6.Lemma.2} Let $T$ be a linear operator with $D(T)^a = X$ and $R(T) \subseteq Y$, where $X$ and $Y$ are $B$-spaces. If $R(T')$ is weakly* dense in $X'$, then $T$ has an inverse. 
\end{lemma}

\begin{proof}
 Suppose that there exists an $x_{0} \neq 0$ such that $Tx_{0} = 0$. Then 
\[
\langle x _{0}, T ^{\prime} y ^{*} \rangle = \langle T x _{0}, y ^{*} \rangle = 0 \quad \text{ for all } \quad y ^{*} \in D (T ^{\prime}). \]

This shows that the weak* closure of $R(T')$ is a proper linear subspace of $X'$, contrary to the hypothesis. 
\end{proof}

\begin{theorem}[R. S. \textsc{Phillips} {\cite{phillips252}}]\label{thm:VIII.6.Theorem.2} Let $X$ be a complex $B$-space, and $T$ a closed linear operator with $D(T)^a = X$ and $R(T) \subseteq X$. Then 
\begin{equation}\label{eq:VIII.6.2}
\varrho (T) = \varrho (T ^{\prime}) \quad \text{ and } \quad R (\lambda ; T) ^{\prime} = R (\lambda ; T ^{\prime}) \quad \text{ for } \quad \lambda \in \varrho (T).
\end{equation}
\end{theorem}

\begin{proof}
 If $\lambda \in \varrho(T)$, then, by Theorem~\ref{thm:VIII.6.Theorem.1}, $\lambda \in \varrho(T')$ and $R(\lambda; T)' = R(\lambda; T')$. On the other hand, if $\lambda \in \varrho(T')$, then Lemma~\ref{lem:VIII.6.Lemma.2} shows that $(\lambda I - T)$ has an inverse $(\lambda I - T)^{-1}$ which is closed with $(\lambda I - T)$. Lemma~\ref{lem:VIII.6.Lemma.1} then shows that $D((\lambda I - T)^{-1}) = R(\lambda I - T)$ is strongly dense in $Y$. Hence, by Theorem~\ref{thm:VIII.6.Theorem.1}, $\lambda \in \varrho(T)$. 
\end{proof}


\yossection{7}{Dunford's Integral}
\label{sec:VIII.7}

Let $X$ be a complex $B$-space and $T$ a bounded linear operator $\in L(X, X)$. We shall define a function $f(T)$ of $T$ by Cauchy's type integral 
\[
f (T) = (2 \pi i) ^{-1} \int_{C} f (\lambda) R (\lambda ; T) \dd \lambda .
\]

To this purpose, we denote by $F(T)$ the family of all complex-valued functions $f(\lambda)$ which are holomorphic in some neighbourhood of the spectrum $\sigma(T)$ of $T$; the neighbourhood need not be connected, and can depend on the function $f(\lambda)$. Let $f \in F(T)$, and let an open set $U \supseteq \sigma(T)$ of the complex plane be contained in the domain of holomorphy of $f$, and suppose further that the boundary $\partial U$ of $U$ consists of a finite number of rectifiable Jordan curves, oriented in positive sense. Then the bounded linear operator $f(T)$ will be defined by 
\begin{equation}\label{eq:VIII.7.1}
f (T) = (2 \pi i) ^{-1} \int_{\partial U} f (\lambda) R (\lambda ; T) \dd \lambda ,
\end{equation}
and the integral on the right may be called a Dunford's integral. By Cauchy's integral theorem, the value $f(T)$ depends only on the function $f$ and the operator $T$, but not on the choice of the domain $U$. 

The following operational calculus holds: 

\begin{theoremn}[N. \textsc{Dunford}] If $f$ and $g$ are in $F(T)$, and $\alpha$ and $\beta$ are complex numbers, then
\begin{gather}
\alpha f+\beta g\text{ is in }F(T)\text{ and }\alpha f(T)+\beta g(T)=(\alpha f+\beta g)(T),\label{eq:VIII.7.2}\\
f\cdot g\text{ is in }F(T)\text{ and }f(T)\cdot g(T)=(f\cdot g)(T),\label{eq:VIII.7.3}
\end{gather}
\begin{equation}\label{eq:VIII.7.4}
\left.\begin{gathered}
\parbox{0.86\linewidth}{\setlength{\parindent}{1.2em}\indent
if $f$ has the Taylor expansion $f(\lambda)=\sum\limits_{n=0}^{\infty}\alpha_n\lambda^n$ convergent in a neighbourhood $U$ of $\sigma(T)$, then}\vphantom{\sum}\mathstrut\\[2pt]
f(T)=\sum\limits_{n=0}^{\infty}\alpha_nT^n\text{ (in the operator norm topology),}
\end{gathered}\right\}
\end{equation}
\begin{equation}\label{eq:VIII.7.5}
\left.\begin{gathered}
\parbox{0.86\linewidth}{\setlength{\parindent}{1.2em}\indent
let $f_n\in F(T)$ $(n=1,2,\ldots)$ be holomorphic in a fixed neighbourhood $U$ of $\sigma(T)$. If $f_n(\lambda)$ converges to $f(\lambda)$ uniformly on $U$, then $f_n(T)$ converges to $f(T)$ in the operator norm topology,}
\end{gathered}\right\}
\end{equation}

\begin{equation}\label{eq:VIII.7.6}
\text{ if } f \in F (T), \text{ then } f \in F (T ^{\prime}) \text{ and } f (T ^{\prime}) = f (T) ^{\prime}.
\end{equation}

\end{theoremn}
\begin{proof}
 \eqref{eq:VIII.7.2} is clear. Proof of \eqref{eq:VIII.7.3}. Let $U_{1}$ and $U_{2}$ be open neighbourhoods of $\sigma(T)$ whose boundaries $\partial U_{1}$ and $\partial U_{2}$ consist of a finite number of rectifiable Jordan curves, and assume that $U_{1} + \partial U_{1} \subseteq U_{2}$ and that $U_{2} + \partial U_{2}$ is contained in the domain of holomorphy of $f$ and $g$. Then, by virtue of the resolvent equation and Cauchy's integral theorem, we obtain 
\[
\begin{aligned}
&f (T) g (T) = - (4 \pi^{2}) ^{-1} \int_ {\partial U _{1}} f (\lambda) R (\lambda ; T) \dd \lambda \cdot \int_ {\partial U _{2}} g (\mu) R (\mu ; T) \dd \mu
\\
&= - (4 \pi^{2}) ^{-1} \int_ {\partial U _{1}} \int_ {\partial U _{2}} f (\lambda) g (\mu) (\mu - \lambda) ^{-1} (R (\lambda ; T) - R (\mu ; T)) \dd \lambda \dd \mu
\\
&= (2 \pi i) ^{-1} \int_ {\partial U _{1}} f (\lambda) R (\lambda ; T) \cdot \left\{(2 \pi i) ^{-1} \int_ {\partial U _{2}} (\mu - \lambda) ^{-1} g (\mu) \dd \mu \right\} \dd \lambda
\\
&- (2 \pi i) ^{-1} \int_ {\partial U _{2}} g (\mu) R (\mu ; T) \cdot \left\{(2 \pi i) ^{-1} \int_ {\partial U _{1}} (\mu - \lambda) ^{-1} f (\lambda) \dd \lambda \right\} \dd \mu
\\
&= (2 \pi i) ^{-1} \int_ {\partial U _{1}} f (\lambda) g (\lambda) R (\lambda ; T) \dd \lambda = (f \cdot g) (T).
\end{aligned}
\]

Proof of \eqref{eq:VIII.7.4}. By hypothesis, $U$ must contain a circle $\{\lambda; |\lambda| \leqq r_{\sigma}(T) + \varepsilon\}$ $\varepsilon > 0$, in its interior, where $r_{\sigma}(T)$ is the spectral radius of $T$ (Theorem~\ref{thm:VIII.2.Theorem.4} in Chapter~\ref{chap:VIII}, \ref{sec:VIII.2}). Hence the power series $f(\lambda) = \sum\limits_{n=0}^{\infty} \alpha_n \lambda^n$ converges uniformly on the circle $C = \{\lambda; |\lambda| \leqq r_{\sigma}(T) + \varepsilon\}$, for some $\varepsilon > 0$. Hence, by Cauchy's integral theorem and Laurent's expansion $R(\lambda; T) = \sum\limits_{n=1}^{\infty} \lambda^{-n} T^{n-1}$ of $R(\lambda; T)$ (Chapter~\ref{chap:VIII}, \ref{sec:VIII.2}, \eqref{eq:VIII.2.5}), 
\[
\begin{aligned}
f (T) &= (2 \pi i) ^{-1} \int_{\partial C} \left(\sum\limits_{k=0} ^{\infty} \alpha_{k} \lambda^{k}\right) R (\lambda ; T) \dd \lambda = (2 \pi i) ^{-1} \sum\limits_{k=0} ^{\infty} \alpha_{k} \int_{\partial C} \lambda^{k} R (\lambda ; T) \dd \lambda
\\
&= (2 \pi i) ^{-1} \sum\limits_{k=0} ^{\infty} \alpha_{k} \sum\limits_{n=1} ^{\infty} \int_{\partial C} \lambda^{k-n} T ^{n-1} \dd \lambda = \sum\limits_{k=0} ^{\infty} \alpha_{k} T ^{k}.
\end{aligned}
\]

\eqref{eq:VIII.7.5} is proved by \eqref{eq:VIII.7.1}, and \eqref{eq:VIII.7.6} is also proved by \eqref{eq:VIII.7.1} and formula \eqref{eq:VIII.7.2} of the preceding section. 
\end{proof}

\begin{corollary}[Spectral Mapping Theorem]\label{cor:VIII.7.Corollary.1} If $f$ is in $F(T)$, then $f(\sigma(T)) = \sigma(f(T))$. 
\end{corollary}

\begin{proof}
 Let $\lambda \in \sigma(T)$, and define the function $g$ by $g(\mu) = (f(\lambda) - f(\mu)) / (\lambda - \mu)$. By the Theorem, $f(\lambda)I - f(T) = (\lambda I - T)g(T)$. Hence, if $(f(\lambda)I - f(T))$ had a bounded inverse $B$, then $g(T)B$ would be the bounded inverse of $(\lambda I - T)$. Thus $\lambda \in \sigma(T)$ implies that $f(\lambda) \in \sigma(f(T))$. Let, conversely, $\lambda \in \sigma(f(T))$, and assume that $\lambda \bar{\in} f(\sigma(T))$. Then the function $g(\mu) = (f(\mu) - \lambda)^{-1}$ must belong to $F(T)$, and so, by the preceding Theorem, $g(T)(f(T) - \lambda I) = I$ which contradicts the assumption $\lambda \in \sigma(f(T))$. 
\end{proof}

\begin{corollary}\label{cor:VIII.7.Corollary.2} If $f \in F(T)$, $g \in F(f(T))$ and $h(\lambda) = g(f(\lambda))$, then $h$ is in $F(T)$ and $h(T) = g(f(T))$. 
\end{corollary}

\begin{proof}
 That $h \in F(T)$ follows from Corollary~\ref{cor:VIII.7.Corollary.1}. Let $U_{1}$ be an open neighbourhood of $\sigma(f(T))$ whose boundary $U_{1}$ consists of a finite number of rectifiable Jordan curves such that $U_{1} + \partial U_{1}$ is contained in the domain of holomorphy of $g$. Let $U_{2}$ be a neighbourhood of $\sigma(T)$ whose boundary $\partial U_{2}$ consists of a finite number of rectifiable Jordan curves such that $U_{2} + \partial U_{2}$ is contained in the domain of holomorphy of $f$ and $f(U_{2} + \partial U_{2}) \subseteq U_{1}$. Then we have, for $\lambda \in \partial U_{1}$, 
\[
R (\lambda ; f (T)) = (2 \pi i) ^{-1} \int_ {\partial U _{2}} (\lambda - f (\mu)) ^{-1} R (\mu ; T) \dd \mu ,
\]
since the right hand operator $S$ satisfies, by \eqref{eq:VIII.7.3}, the equation $(\lambda I - f(T)) S = S(\lambda I - f(T)) = I$. Hence, by Cauchy's integral theorem, 
\[
\begin{aligned}
g (f (T)) &= (2 \pi i) ^{-1} \int_ {\partial U _{1}} g (\lambda) R (\lambda ; f (T)) \dd \lambda
\\
\qquad &= (- 4 \pi^{2}) ^{-1} \int_ {\partial U _{1}} \int_ {\partial U _{2}} g (\lambda) (\lambda - f (\mu)) ^{-1} R (\mu ; T) \dd \mu \dd \lambda
\\
\qquad &= (2 \pi i) ^{-1} \int_ {\partial U _{2}} R (\mu ; T) g (f (\mu)) \dd \mu = h (T). \end{aligned}
\]
\end{proof}

\begin{remarku} The introduction of the operational calculus based on a formula like \eqref{eq:VIII.7.1} goes back to the investigations by H. \textsc{Poincare} on  
continuous groups ($1899$). The exposition of the operational calculus in this section is adapted from N. \textsc{Dunford}-J. \textsc{Schwartz} \cite{dunford44}. In the next chapter on semi-groups, we shall frequently make use of Dunford's integral for a closed unbounded operator $T$. 
\end{remarku}


\yossection{8}{The Isolated Singularities of a Resolvent}
\label{sec:VIII.8}

Let $\lambda_{0}$ be an isolated singular point of the resolvent $R(\lambda; T)$ of a closed linear operator $T$ on a complex $B$-space $X$ into $X$. Then $R(\lambda; T)$ can be expanded into Laurent series 
\begin{equation}\label{eq:VIII.8.1}
R (\lambda ; T) = \sum\limits_{n = - \infty} ^{\infty} (\lambda - \lambda_{0}) ^{n} A _{n}, A _{n} = (2 \pi i) ^{-1} \int_{C} (\lambda - \lambda_{0}) ^{-n-1} R (\lambda ; T) \dd \lambda ,
\end{equation}
where $C$ is a circumference of sufficiently small radius: $|\lambda - \lambda_{0}| = \varepsilon$ such that the circle $|\lambda - \lambda_{0}| \leqq \varepsilon$ does not contain other singularities than $\lambda = \lambda_{0}$, and the integration is performed counter-clockwise. By virtue of the resolvent equation we obtain 

\begin{theorem}\label{thm:VIII.8.Theorem.1} $A$'s are mutually commutative bounded linear operators and 
\begin{equation}\label{eq:VIII.8.2}
\begin{gathered}
T A _{k} x = A _{k} T x \text{ for } x \in D (T) (k = 0, \pm 1, \pm 2, \dots), A _{k} A _{m} = 0 \quad \text{ for } \quad k \geqq 0, m \leqq - 1,
\\
A _{n} = (- 1) ^{n} A _{0} ^{n+1} (n \geqq 1),
\\
A _{-p-q+1} = A _{-p} A _{-q} (p, q \geqq 1).
\end{gathered}
\end{equation}
\end{theorem}

\begin{proof}
 The boundedness and the mutual commutativity of $A$'s and the commutativity of $A$'s with $T$ are clear from the integral representation of $A$'s. 


We substitute the expansion of $R(\lambda; T)$ in the resolvent equation $R(\lambda; T) - R(\mu; T) = (\mu - \lambda) R(\lambda; T) R(\mu; T)$, and obtain 
\[
\sum\limits_{k = - \infty} ^{\infty} A _{k} \frac {(\lambda - \lambda_{0}) ^{k} - (\mu - \lambda_{0}) ^{k}}{(\lambda - \lambda_{0}) - (\mu - \lambda_{0})} = - \sum\limits_{k, m = - \infty} ^{\infty} A _{k} A _{m} (\lambda - \lambda_{0}) ^{k} (\mu - \lambda_{0}) ^{m}.
\]

The coefficient of $A_{n}$ on the left is 
\[
\begin{aligned}
(\lambda - \lambda_{0}) ^{n-1} + (\lambda - \lambda_{0}) ^{n-2} (\mu - \lambda_{0}) + \dots + (\mu - \lambda_{0}) ^{n-1}, n& \geqq 1,
\\
& - \{(\lambda - \lambda_{0}) ^{n} (\mu - \lambda_{0}) ^{-1}
\\
&\quad + (\lambda - \lambda_{0}) ^{n+1} (\mu - \lambda_{0}) ^{-2} + \dots
\\
+ (\lambda - \lambda_{0}) ^{-1} (\mu - \lambda_{0}) ^{n} \}, n& < 0.
\end{aligned}
\]

Hence the terms containing $(\lambda - \lambda_0)^k (\mu - \lambda_0)^m$ with $k \geqq 0$ and $m \leqq -1$ are missing so that we must have the orthogonality $A_k A_m = 0$ ( $k \geqq 0$, $m \leqq -1$). Hence 
\[
R ^{+} (\lambda ; T) = \sum\limits_{n=0} ^{\infty} A _{n} (\lambda - \lambda_{0}) ^{n} \quad \text{ and } \quad R ^{-} (\lambda ; T) = \sum\limits_{n = - \infty} ^{-1} A _{n} (\lambda - \lambda_{0}) ^{n}
\]
must both satisfy the resolvent equation. Substituting the expansion of $R^{+}(\lambda; T)$ in the resolvent equation 
\[
R ^{+} (\lambda ; T) - R ^{+} (\lambda ; T) = (\mu - \lambda) R ^{+} (\lambda ; T) R ^{+} (\mu ; T),
\]
we obtain, setting $(\lambda - \lambda_0) = h$, $(\mu - \lambda_0) = k$, 
\[
\sum\limits_{p=1} ^{\infty} A _{p} \left(h ^{p} - k ^{p}\right) = (k - h) \left(\sum\limits_{p=0} ^{\infty} A _{p} h ^{p}\right) \left(\sum\limits_{q=0} ^{\infty} A _{q} k ^{q}\right).
\]

Hence, dividing both sides by $(k - h)$, we obtain 
\[
- \sum\limits_{p=1} ^{\infty} A _{p} (h ^{p-1} + h ^{p-2} k + \dots + k ^{p-1}) = \sum\limits_{p,q=0} ^{\infty} h ^{p} k ^{q} A _{p} A _{q}
\]
so that we have $-A_{p+q+1} = A_pA_q$ ( $p, q \geqq 0$). Thus, in particular, 
\[
A _{1} = - A _{0} ^{2}, A _{2} = - A _{1} A _{0} = (- 1) ^{2} A _{0} ^{3}, \dots , A _{n} = (- 1) ^{n} A _{0} ^{n+1} (n \geqq 1).
\]

Similarly, from the resolvent equation for $R^{-}(\lambda; T)$, we obtain, setting $(\lambda - \lambda_{0})^{-1} = h$, $(\mu - \lambda_{0})^{-1} = k$, 
\[
\sum\limits_{p=1} ^{\infty} A _{-p} \left(h ^{p-1} + h ^{p-2} k + \dots + k ^{p-1}\right) = \sum\limits_{p,q=1} ^{\infty} h ^{p-1} k ^{q-1} A _{-p} A _{-q},
\]
so that we have $A_{-p-q+1} = A_{-p}A_{-q}$ ( $p, q \geqq 1$). In particular, we have 
\[
A _{-1} = A _{-1} ^{2}, A _{-2} = A _{-1} A _{-2}, \dots , A _{-n} = A _{-1} A _{-n} (n \geqq 1).
\]
\end{proof}

\begin{theorem}\label{thm:VIII.8.Theorem.2} We have 
\begin{equation}\label{eq:VIII.8.3}
\begin{gathered}
{A _{n} = (T - \lambda_{0} I) A _{n+1} (n \geqq 0),} {(T - \lambda_{0} I) A _{-n} = A _{-(n+1)} = (T - \lambda_{0} I) ^{n} A _{-1} \quad (n \geqq 1),}
\\
{(T - \lambda_{0} I) A _{0} = A _{-1} - I.}
\end{gathered}
\end{equation}
\end{theorem}

\begin{proof}
 By the integral representation of $A_{k}$, we see that the range $R(A_{k})$ is in the domain of $T$, so that we can multiply $A_{k}$ by $T$ on the left. Thus our Theorem is proved by the identity 
\begin{equation}\label{eq:VIII.8.4}
\begin{aligned}
&I = (\lambda I - T) \sum\limits_{k = - \infty} ^{\infty} A _{k} (\lambda - \lambda_{0}) ^{k} = \{(\lambda - \lambda_{0}) I + (\lambda_{0} I - T) \} \sum\limits_{k = - \infty} ^{\infty} A _{k} (\lambda - \lambda_{0}) ^{k}.
\end{aligned}
\end{equation}
\end{proof}

\begin{theorem}\label{thm:VIII.8.Theorem.3} If $\lambda_0$ is a pole of $R(\lambda; T)$ of order $m$, then $\lambda_0$ is an eigenvalue of $T$. We have 

$R(A_{-1}) = N((\lambda_0I - T)^n)$ and $R(I - A_{-1}) = R((\lambda_0I - T)^n)$ for $n \geqq m$, \eqref{eq:VIII.8.4} so that, in particular, 
\begin{equation}\label{eq:VIII.8.5}
X = N ((\lambda_{0} I - T) ^{n}) \oplus R ((\lambda_{0} I - T) ^{n}) \quad \text{ for } \quad n \geqq m.
\end{equation}
\end{theorem}

\begin{proof}
 Since $A_{-1}$ is a bounded linear operator satisfying $A_{-1}^2 = A_{-1}$, it is easy to see that 
\begin{equation}\label{eq:VIII.8.6}
N (A _{-1}) = R (I - A _{-1}). \end{equation}

We put $X_{1} = N(A_{-1}) = R(I - A_{-1})$, and put also 
\begin{equation}\label{eq:VIII.8.7}
X _{2} = R (A _{-1}), N _{n} = N ((\lambda_{0} I - T) ^{n}) \text{ and } R _{n} = R _{n} ((\lambda_{0} I - T) ^{n})
\end{equation}

Let $x \in N_n$ where $n \geqq 1$. Then we see, by $(T - \lambda_0 I)^n A_{n-1} = (T - \lambda_0 I) A_0 = A_{-1} - I$, that $0 = A_{n-1}(T - \lambda_0 I)^n x = (T - \lambda_0 I)^n A_{n-1} x = (T - \lambda_0 I) A_0 x = A_{-1} x - x$ so that $x = A_{-1} x \in X_2$. Thus $N_n$ with $n \geqq 1$ belongs to $X_2$. Let, conversely, $x \in X_2$. Then we have $x = A_{-1} y$ and so $x = A_{-1} A_{-1} y = A_{-1} x$ by $A_{-1} = A_{-1}^2$; consequently, we have $(T - \lambda_0 I)^n x = A_{-(n+1)} x$ by $(T - \lambda_0 I)^n A_{-1} = A_{-(n+1)}$. Since $A_{-(n+1)} = 0$ for $n \geqq m$ by hypothesis, it follows that $X_2 \subseteq N_n$ for $n \geqq m$ and so 
\begin{equation}\label{eq:VIII.8.8}
N _{n} = X _{2} \quad \text{ if } \quad n \geqq m.
\end{equation}

Because $(T - \lambda_0I)A_{-m} = A_{-(m+1)} = 0$ and $A_{-m}\neq 0$, the number $\lambda_0$ is an eigenvalue of $T$. 


We see that $X_{1} = N(A_{-1}) = R(I - A_{-1}) \subseteq R_{n}$ by $(T - \lambda_0 I)^n A_{n-1} = A_{-1} - I$. If $n \geqq m$, then $x \in R_n \cap N_n$ implies $x = 0$. For, if $x = (\lambda_0 I - T)^n y$ and $(\lambda_0 I - T)^n x = 0$, then, by \eqref{eq:VIII.8.8}, $y \in N_{2n} = N_n$ and therefore $x = 0$. Next suppose $x \in R_n$ with $n \geqq m$, and write $x = x_1 + x_2$ where $x_1 = (I - A_{-1})x \in X_1$, $x_2 = A_{-1}x \in X_2$. Then, since $X_1 \subseteq R_n$, $x_2 = x - x_1 \in R_n$. But $x_2 \in X_2 = N_n$ by \eqref{eq:VIII.8.8}, and so $x_2 \in R_n \cap N_n$, $x_2 = 0$. This proves that $x = x_1 \in X_1$. Therefore we have proved that $R_n = X_1$ if $n \geqq m$.
\end{proof}

\begin{theorem}\label{thm:VIII.8.Theorem.4} If, in particular, $T$ is a bounded linear operator such that $X_{2} = R(A_{-1})$ is a finite dimensional linear subspace of $X$, then $\lambda_{0}$ is a pole of $R(\lambda; T)$. 
\end{theorem}

\begin{proof}
 Let $x_1, x_2, \ldots, x_k$ be a base of the linear space $X_2$. Since $x_1, Tx_1, T^2x_1, \ldots, T^k x_1$ are linearly dependent vectors of $X_2$, there exists a non-zero polynomial $P_1(\lambda)$ such that $P_1(T)x_1 = 0$. Similarly there exist non-zero polynomials $P_2(\lambda), \ldots, P_k(\lambda)$ such that $P_j(T)x_j = 0$ ( $j = 2, 3, \ldots, k$). Then for the polynomial $P(\lambda) = \prod\limits_{j=1}^{k} P_j(\lambda)$, we must have $P(T)x_j = 0$ ( $j = 1, \ldots, k$) and hence $P(T)x = 0$ for every $x \in X_2$. Let 
\[
P (\lambda) = \alpha \prod\limits_{j=0} ^{s} (\lambda - \lambda_{j}) ^ {\nu_{j}} (\alpha \neq 0)
\]
be the factorization of $P(\lambda)$. Then we can prove that $(T - \lambda_0 I)^{\nu_0} x = 0$ for every $x \in X_2$. Assume the contrary, and let $x_0 \in X_2$ be such that $(T - \lambda_0 I)^{\nu_0} x_0 \neq 0$. Then, by $P(T)x_0 = 0$, we see that there exist at least one $\lambda_j (j \neq 0)$ and a polynomial $Q(\lambda)$ such that 
\[
(T - \lambda_{j} I) Q (T) (T - \lambda_{0} I) ^ {\nu_{0}} x _{0} = 0
\]
and $y = Q(T)(T - \lambda_0 I)^{\nu_0} x_0 \neq 0$. Thus $y \in X_2$ is an eigenvector of $T$ corresponding to the eigenvalue $\lambda_j$. Hence $(\lambda I - T)y = (\lambda - \lambda_j)y$ and so, multiplying both sides by $R(\lambda; T)$, we obtain $y = (\lambda - \lambda_j)R(\lambda; T)y$ which implies 
\[
y = A _{-1} y = (2 \pi i) ^{-1} \int_{C} R (\lambda ; T) y \dd \lambda = (2 \pi i) ^{-1} \int_{C} (\lambda - \lambda_{j}) ^{-1} y \dd \lambda = 0,
\]
by taking the circumference $C$ with $\lambda_{0}$ as centre sufficiently small. This is a contradiction, and so there must exist a positive integer $m$ such that $(T-\lambda_{0}I)^{m}X_{2}=0$. Thus, by $X_{2}=R(A_{-1})$ and $(T-\lambda_{0}I)^{n}A_{-1}=A_{-(n+1)}$, we see that $A_{-(n+1)}=0$ for $n\geqq m$. 
\end{proof}

\yossub{Comments and References}

Section~\hyperref[sec:VIII.6]{$6$} is adapted from R. S. \textsc{Phillips} \cite{phillips252}. Section~\hyperref[sec:VIII.8]{$8$} is adapted from M. \textsc{Nagumo} \cite{nagumo215} and A. \textsc{Taylor} \cite{taylor321}. Parts of these sections can easily be extended to the case of a locally convex linear topological space. See, e.g., section $13$ of the following chapter. 
